\section{Limitation and future work}
\label{append:limit}

\subsection{Limitations of our proposed approaches}

\paragraph{Approach 1}
The Approach 1 provides a broader solution by predicting the allocation of computational resources across non-vocabulary parameters, vocabulary parameters, and training data based on experimental data points. This method's strength lies in its holistic view, allowing for a balanced resource distribution that potentially enhances model efficiency and performance. However, this approach is constrained by the granularity and range of the experimental data points available, which can introduce errors in the fitting process. The requirement for substantial computational resources to perform these fittings may also limit its accessibility and scalability. Despite these challenges, when experimental data is ample and computational resources are sufficient, the Approach 1 can significantly refine the precision of resource allocation decisions in the development of large-scale language models.

\paragraph{Approach 2} 
By calculating the derivative of FLOPs with respect to the vocabulary size and solving for zero, this approach fundamentally relies on the precision of the FLOPs equation and our tokens-characters relationship function. Further, this method does not allow us to independently determine the optimal allocation of non-vocabulary parameters and training data size. Therefore, it necessitates information about the relationships between these attributes and the FLOPs budget from the experimentally fitted scaling laws, making this approach less useful in practice. Despite these limitations, the derivative-based approach offers notable advantages, including closely matched predictions with the scaling laws derived from actual experimental data in the Approach 2. Furthermore, its reliance on numerical solutions rather than exhaustive deep learning experiments makes it rapid and broadly applicable across various tokenizers, highlighting its utility in preliminary model configuration stages where quick estimates are key.

\paragraph{Approach 3}
Similar with the Approach 1, the proposed Approach 3 requires multiple experimental runs across different non-vocabulary parameters, vocabulary sizes and number of training data. Therefore, the approach is constrained by the granularity and range of the experimental data points available to some extent. However, the proposed Approach 3 is flexible that it considers the fact that the non-vocabulary parameters and the number of training data are not always following the compute-optimal scaling laws \citep{hoffmann2022training}, \textit{i.e.}, equal scaling, in real-world applications.

\subsection{Larger models and different architectures}

We have shown that our predictions hold for models with up to three billion parameters (\autoref{sec:discussion}). However, LLMs are often orders of magnitude larger, such as the 400-billion parameter Llama-3 model \cite{metallama3blog}. Further, we have decided to focus on dense transformer language models, as they are most commonly used for LLMs. However, many non-transformer models have been proposed and scaled up to billions of parameters \cite{peng2023rwkv,peng2024eagle}. Exploring to what extent our findings hold in even larger models and with different architectures is a promising direction for future work.

\subsection{Parametric function for the loss when considering the vocabulary}
Researchers~\citep{hoffmann2022training,muennighoff2024scaling} consider modeling the language modeling loss with parametric functions in the form of $\mathcal{L}=P_1+P_2/N^{\alpha}+P_3/D^{\beta}$, where $\{P_1,P_2,P_3,\alpha,\beta\}$ are learnable variables. The first term of loss represents the minimum achievable loss, and the second and third terms represent the contribution to the loss from the model size $N$ and number of training tokens $D$. The parametric function allows predicting the loss $L$ given $N$ and $D$ even if ($N$,$D$) are not optimally allocated. In prior work, this loss formula accounts for changes in model size and training data but does not explicitly address the complexities introduced by varying vocabulary sizes. Incorporating vocabulary size into the loss predictor is challenging: Vocabulary size affects the model directly as well as the number of training tokens and the quality of tokenization by the tokenizer. A tokenizer with a large vocabulary size makes it easier to capture semantic information in raw text and reduces the frequency of out-of-vocabulary words. For instance, a large vocabulary size may cover common phrases, common subwords, and specialized terminology. Therefore, even if the same number of tokens are trained, the performance of the model trained on tokens with different qualities will be different.

Future work in this area could explore various parametric non-linear loss functions to predict the interactions between vocabulary size, model size, and training data with different compute allocations, not just the case of optimal compute allocation. Additionally, empirical studies on different datasets could help in understanding how vocabulary size impacts loss under varied data conditions, guiding the development of more adaptive loss prediction models.

\subsection{Extensions to multilingual and multimodal scenarios}
Future work could extend the proposed approaches to encompass multilingual and multimodal scenarios. Multilingual models require a nuanced understanding of vocabulary due to linguistic diversity, which may affect the optimal vocabulary size and the computation of FLOPs differently across languages. Adapting these methods to consider linguistic features and tokenization variations could lead to more tailored and efficient resource allocations for multilingual models. Different languages compete with each other for the model's ability to allocate to that language \citep{blevins2024breaking}, which makes it necessary to take into account the relationship between different languages when setting the size of word lists for different languages in a multilingual scenario.

For multimodal models that integrate text with other data types such as images or video, the optimal vocabulary size might interact uniquely with non-linguistic parameters. Recent work \citep{aghajanyan2022cm3,tian2024visual} models visual concepts in an autoregressive manner with tokenization like the processing of text data. It is interesting to explore the size of visual vocabulary size, \textit{i.e.}, the codebook size \citep{esser2021taming}, in the visual tasks and vision-language tasks. How to set the vocabulary size and the compute resource efficiently for different modalities remains an open issue.

\section{Potential social impact}
\label{append:social}

The positive potential social impact of this research on vocabulary size in language model scaling is substantial. By optimizing large language models with the consideration of the vocabulary size and other attributes jointly, the paper provides a foundational understanding that can lead to more lightweight and cost-effective pre-trained large language models. This efficiency can democratize access to advanced language processing technologies, making it feasible for smaller organizations and the general public to benefit from powerful AI tools. Such advancements can benefit various domains, for example, improve accessibility features for individuals with disabilities, where efficient language models can be used to analyze medical records and assist in diagnostics. Furthermore, the reduction in computational requirements for training these models can lead to a decrease in energy usage, contributing positively to environmental sustainability efforts.

On the other hand, the misuse of pretrained language models may pose risks, including the creation of highly realistic deepfakes that can spread disinformation and undermine trust in media and institutions. These models can generate misleading content, automate cyberattacks through convincing phishing schemes, and produce large-scale spam, degrading online communication. Additionally, they can be used to generate harmful or abusive content, such as hate speech, which perpetuates discrimination and harms vulnerable populations. To mitigate these risks, it is crucial to develop trustworthy language models, implement robust monitoring systems, and foster collaboration among researchers, policymakers, and users.

